\section{Método de prototipado}

El prototipo se construyó en \textbf{Figma} a partir de la arquitectura de información, el
mapa de navegación y los hallazgos documentados en la Actividad 3. La herramienta permitió
trabajar con componentes reutilizables, variables, estilos compartidos y conexiones entre
pantallas.

\subsection{Decisión de fidelidad}

La alta fidelidad no se trató como una colección de imágenes aisladas. Cada flujo incorporó
contenido representativo, jerarquía visual, controles con estados, objetivos de tarea y una
secuencia navegable. Este enfoque coincide con la función que Hartson y Pyla (2019) asignan al
prototipo de alta fidelidad: aproximar la experiencia del producto para examinar tanto su
apariencia como su operación.

\subsection{Cobertura de los cuatro requisitos}

\begin{uxlist}
  \lead{Representación fiel}{La paleta, la tipografía, la retícula y los componentes se aplicaron de manera consistente en los cinco flujos.}
  \lead{Interacciones evaluables}{Las 30 pantallas quedaron vinculadas mediante 84 conexiones sin destinos inválidos.}
  \lead{Comentarios visuales detallados}{La separación entre fundamentos, componentes y prototipos facilita revisar decisiones concretas de estética y legibilidad.}
  \lead{Comprensión del producto final}{Los flujos cubren consulta, análisis, riesgo, publicación, permisos y navegación integral del portal.}
\end{uxlist}

\subsection{Acceso al archivo de trabajo}

El prototipo y la guía se encuentran en el archivo duplicado
\href{https://www.figma.com/design/4RA9Q606HpPNveRruoHr7t/Karisma-Data---Actividad-4--Dev?node-id=0-1}{Karisma Data - Actividad 4 - Dev}.
La copia conservó intacta la versión de origen y concentró los ajustes de esta actividad.
